% Zone „Das kennst du schon“ – aus bank/reelle-zahlen/zone.jsonl (zusammenbau v0.1)
\einheitenkopf[zone][Kennst du schon]{Das kennst du schon}
\setcounter{aufgabe}{0}

%% TODO Titel: Ich-kann-Satz fehlt in der Bank (Platzhalter: Kettenname) – Nr. 1
%% TODO Anweisung: ein Satz über der ersten Teilaufgabe fehlt in der Bank – Nr. 1
\begin{aufgabe}{Bruch in Dezimalzahl umwandeln (Division), abbrechende und periodische Dezimalzahlen erkennen und mit Periodenstrich schreiben, Dezimalzahl in Bruch}
\begin{teile}
\teil $\frac{3}{4}$ – als Dezimalzahl? \leerfeld
\teil $\frac{7}{20}$ – als Dezimalzahl? \leerfeld
\end{teile}
\end{aufgabe}

%% TODO Titel: Ich-kann-Satz fehlt in der Bank (Platzhalter: Kettenname) – Nr. 2
%% TODO Anweisung: ein Satz über der ersten Teilaufgabe fehlt in der Bank – Nr. 2
\begin{aufgabe}{Zahlengerade mit negativen Zahlen, Brüche und Dezimalzahlen eintragen, Zahlen vergleichen und ordnen}
\begin{teile}
\teil $-2$ oder $-6$ – welche Zahl ist größer? \leerfeld
\teil $0{,}8$; $-0{,}9$; $0{,}85$ – aufsteigend geordnet? \leerfeld ; \leerfeld ; \leerfeld
\end{teile}
\end{aufgabe}

%% TODO Titel: Ich-kann-Satz fehlt in der Bank (Platzhalter: Kettenname) – Nr. 3
%% TODO Anweisung: ein Satz über der ersten Teilaufgabe fehlt in der Bank – Nr. 3
\begin{aufgabe}{Quadratwurzel als Umkehrung des Quadrierens, Quadratzahlen bis zwanzig hoch zwei, Wurzel mit dem Taschenrechner und Näherungswert, Wurzel zwischen zwei Nachbar-Quadratzahlen abschätzen, keine Wurzel aus negativer Zahl}
\begin{teile}
\teil $\sqrt{81}$ – welche Zahl? \leerfeld
\teil $\sqrt{30}$ – zwischen welchen zwei ganzen Zahlen? zwischen \leerfeld und \leerfeld
\end{teile}
\end{aufgabe}

%% TODO Titel: Ich-kann-Satz fehlt in der Bank (Platzhalter: Kettenname) – Nr. 4
%% TODO Anweisung: ein Satz über der ersten Teilaufgabe fehlt in der Bank – Nr. 4
\begin{aufgabe}{Taschenrechner}
\begin{teile}
\teil $\sqrt{7}$ mit dem Taschenrechner – auf zwei Stellen gerundet? \leerfeld
\teil $1{,}3^2$ mit dem Taschenrechner – Anzeige? \leerfeld
\end{teile}
\end{aufgabe}

%% TODO Titel: Ich-kann-Satz fehlt in der Bank (Platzhalter: Kettenname) – Nr. 5
%% TODO Anweisung: ein Satz über der ersten Teilaufgabe fehlt in der Bank – Nr. 5
\begin{aufgabe}{Terme zusammenfassen}
\begin{teile}
\teil $4a + 3a$ – zusammengefasst? \leerfeld
\teil $2{,}5y + 1{,}5y - y$ – zusammengefasst? \leerfeld
\end{teile}
\end{aufgabe}

%% TODO Titel: Ich-kann-Satz fehlt in der Bank (Platzhalter: Kettenname) – Nr. 6
%% TODO Anweisung: ein Satz über der ersten Teilaufgabe fehlt in der Bank – Nr. 6
\begin{aufgabe}{Potenz als Malkette mit Basis und Exponent, Potenz ausrechnen (auch negative Basis), a hoch null gleich eins und negative Hochzahl als Bruch}
\begin{teile}
\teil $3^3$ – als Malkette und ausgerechnet? \leerfeld
\teil $(-2)^3$ – ausgerechnet? \leerfeld
\end{teile}
\end{aufgabe}

%% TODO Titel: Ich-kann-Satz fehlt in der Bank (Platzhalter: Kettenname) – Nr. 7
%% TODO Anweisung: ein Satz über der ersten Teilaufgabe fehlt in der Bank – Nr. 7
\begin{aufgabe}{Brüche als Exponenten lesen und mit Brüchen rechnen}
\begin{teile}
\teil $\frac{1}{2} + \frac{1}{2}$ – ausgerechnet? \leerfeld
\teil $\frac{1}{2} + \frac{1}{3}$ – ausgerechnet? \leerfeld
\end{teile}
\end{aufgabe}

%% TODO Titel: Ich-kann-Satz fehlt in der Bank (Platzhalter: Kettenname) – Nr. 8
%% TODO Anweisung: ein Satz über der ersten Teilaufgabe fehlt in der Bank – Nr. 8
\begin{aufgabe}{Potenz als Malkette mit Basis und Exponent, Potenz ausrechnen (auch negative Basis), a hoch null gleich eins und negative Hochzahl als Bruch – Fehler finden}
\begin{teile}
\teil Finde den Fehler und rechne richtig: \rechnung{3^{-3} &= -27}
\end{teile}
\end{aufgabe}

%% TODO Titel: Ich-kann-Satz fehlt in der Bank (Platzhalter: Kettenname) – Nr. 9
%% TODO Anweisung: ein Satz über der ersten Teilaufgabe fehlt in der Bank – Nr. 9
\begin{aufgabe}{Potenz als Malkette mit Basis und Exponent, Potenz ausrechnen (auch negative Basis), a hoch null gleich eins und negative Hochzahl als Bruch – selbst rechnen}
\begin{teile}
\teil $2^{-4}$ – als Bruch? \leerfeld
\end{teile}
\end{aufgabe}
